\chapter{Presupuesto}

El presente capítulo estima el costo de implementación del prototipo propuesto. El análisis
considera los componentes electrónicos principales, los materiales de fabricación y los
elementos de ensamblaje necesarios para construir una estación inteligente de control ESD
de bajo costo. Los precios utilizados corresponden a valores de referencia obtenidos de
proveedores comerciales en línea, por lo que pueden variar según disponibilidad, impuestos,
envío y fecha de compra.

\section{Estimación de costos de hardware}

La Tabla~\ref{tab:presupuesto_hardware} presenta los principales módulos electrónicos
seleccionados para la implementación del prototipo. Estos componentes corresponden a los
subsistemas de verificación ESD, identificación de usuario, interfaz gráfica, monitoreo
ambiental, almacenamiento y alimentación del sistema.

\begin{table}[H]
\centering
\caption{Estimación de costos de los componentes electrónicos principales.}
\label{tab:presupuesto_hardware}
\begin{tabular}{p{5.2cm} c c c}
\toprule
\textbf{Componente} & \textbf{Cantidad} & \textbf{Costo unitario (USD)} & \textbf{Subtotal (USD)} \\
\midrule
Probador HZR-171 para pulsera antiestática~\cite{HZR171} & 1 & 29.99 & 29.99 \\
ESP32-WROOM-32~\cite{ESP32WROOM32Amazon} & 1 & 9.39 & 9.39 \\
Pantalla HMI ESP32 de 7 pulgadas~\cite{ElecrowESP32HMI7} & 1 & 49.99 & 49.99 \\
Lector de código de barras GM67~\cite{GM67Ebay} & 1 & 32.62 & 32.62 \\
Módulo ESP32-S3 CAM~\cite{FreenoveESP32S3CAM} & 1 & 21.95 & 21.95 \\
Módulo SHT31-D de temperatura y humedad~\cite{SHT31Module} & 1 & 4.7 & 4.7 \\
Módulo USB-C PD~\cite{AITRIPPDTrigger} & 1 & 5.00 & 5.00 \\
Módulo convertidor DC-DC XL6019~\cite{XL6019BoostModule} & 1 & 3.00 & 3.00 \\
Parlante de 8 $\Omega$ 3W \cite{Speaker8Ohm3W} & 1 & 4.00 & 4.00 \\
Memoria microSD de 8 GB \cite{MicroSD8GB} & 1 & 9.99 & 9.99 \\
\midrule
\textbf{Subtotal de hardware} &  &  & \textbf{170.63} \\
\bottomrule
\end{tabular}
\end{table}

La memoria microSD de 8 GB se selecciona para almacenar los registros históricos y
reportes generados por el sistema. Aunque los archivos de reporte en formato CSV requieren
poco espacio de almacenamiento, se selecciona esta capacidad debido a su bajo costo y a que
permite mantener registros durante periodos prolongados sin comprometer la capacidad de
memoria del prototipo. De forma complementaria, el módulo ESP32-S3 CAM dispone de
almacenamiento local que puede utilizarse para guardar información asociada al subsistema
de reconocimiento facial, como imágenes de entrenamiento o vectores de identificación por lo
que no se requiere de una unidad de almacenamiento adicional.

\section{Costos de fabricación y ensamblaje}

Además de los módulos electrónicos principales, se consideran materiales de fabricación
para la interconexión, montaje y construcción física del prototipo. Estos elementos se agrupan
en categorías generales debido a que corresponden a materiales de bajo costo, tales como
cables, conectores, placas de prototipado, tornillería y consumibles de impresión 3D.

\begin{table}[H]
\centering
\caption{Estimación de costos de fabricación y ensamblaje.}
\label{tab:presupuesto_fabricacion}

\begin{tabularx}{\textwidth}{X X c}
\toprule
\textbf{Concepto} &
\textbf{Descripción} &
\textbf{Costo estimado (USD)} \\
\midrule

Cableado eléctrico &
Cable 24 AWG para interconexión interna \cite{AWG24Wire} &
9.66 \\

Placas FR4 de prototipado &
Placas FR4 de 10 cm $\times$ 15 cm \cite{FR4ProtoBoard} &
14.29 \\

Material de impresión 3D &
Consumo aproximado de 500 g de PLA \cite{PLA500gCost} &
7.25 \\

Materiales mecánicos menores &
Tornillos, separadores, conectores, termocontraíble y otros consumibles &
20.00 \\

\midrule
\textbf{Subtotal de fabricación y ensamblaje} &
&
\textbf{51.20} \\

\bottomrule
\end{tabularx}
\end{table}

No se incluye el costo de herramientas o equipos ya disponibles, tales como impresora 3D,
estación de soldadura, computadora de desarrollo, multímetro, instrumentos de medición o
licencias de software. Estos recursos forman parte de la infraestructura disponible para la
ejecución del proyecto y no corresponden a materiales consumibles del prototipo.

\section{Costo total estimado del prototipo}

A partir de los costos de hardware, fabricación y ensamblaje, se obtiene el costo total
estimado mostrado en la Tabla~\ref{tab:presupuesto_total}. Este valor representa el costo
directo de materiales necesarios para la construcción del prototipo funcional.

\begin{table}[H]
\centering
\caption{Costo total estimado del prototipo.}
\label{tab:presupuesto_total}
\begin{tabular}{p{8cm} c}
\toprule
\textbf{Categoría} & \textbf{Costo estimado (USD)} \\
\midrule
Componentes electrónicos principales & 170.63 \\
Materiales de fabricación y ensamblaje & 51.20 \\
\midrule
\textbf{Costo total estimado} & \textbf{221.83} \\
\bottomrule
\end{tabular}
\end{table}

Considerando posibles variaciones en precios, disponibilidad, envío y materiales menores
no contemplados individualmente, se establece un costo de referencia aproximado de
\textbf{250 USD} para la implementación del prototipo.

\section{Comparación con soluciones comerciales}

Uno de los objetivos principales del proyecto es desarrollar una solución de bajo costo en
comparación con sistemas comerciales de control ESD. La Tabla~\ref{tab:comparacion_costos}
presenta una comparación entre el costo estimado del prototipo y diferentes soluciones
comerciales analizadas en el estado del arte.

\begin{table}[H]
\centering
\caption{Comparación económica entre el prototipo propuesto y soluciones comerciales.}
\label{tab:comparacion_costos}
\begin{tabular}{p{7.2cm} c c}
\toprule
\textbf{Sistema} & \textbf{Costo aproximado (USD)} & \textbf{Referencia} \\
\midrule
Prototipo propuesto & 250 & Elaboración propia \\
SCS 770758 Dual Combination Tester & 1181 & \cite{SCS_770758} \\
PGT120 Combination ESD Tester & 2137 & \cite{PGT120} \\
ESDMAN ESD Access Control System & 2500 & \cite{ESDMAN_0016809} \\
SmartLog Pro 2 & 5643 & \cite{SmartLogPro2} \\
PGT130.DT Data Logging ESD Tester & 7676 & \cite{PGT130DT} \\
\bottomrule
\end{tabular}
\end{table}

Tomando como referencia el costo aproximado de 250 USD, el prototipo propuesto presenta una ventaja económica significativa 
frente a los sistemas comerciales avanzados analizados. En particular, el sistema SCS 770758 posee un costo aproximadamente 
4.7 veces superior al del prototipo, mientras que el PGT120 alcanza un costo 8.5 veces mayor. 

La diferencia económica resulta aún más significativa al comparar el prototipo con plataformas industriales más avanzadas. 
El sistema ESDMAN presenta un costo cercano a 10 veces el valor de la solución propuesta, por su parte el 
SmartLog Pro 2 posee un costo aproximado 22.6 veces superior, mientras que 
el PGT130.DT alcanza un valor cercano a 30.7 veces el costo estimado del prototipo.


\section{Análisis económico}

El presupuesto estimado evidencia que es factible implementar una estación 
inteligente de control ESD con un costo aproximado de 250 USD. Esta característica 
resulta especialmente relevante para laboratorios académicos, centros de investigación 
y organizaciones que requieren mejorar la trazabilidad de procedimientos ESD sin realizar 
inversiones de varios miles de dólares en sistemas industriales especializados.

Considerando que algunas de las soluciones comerciales analizadas presentan costos 
entre 5 y más de 30 veces superiores al costo del prototipo propuesto, la arquitectura 
desarrollada ofrece una alternativa económicamente atractiva para aplicaciones de investigación, 
validación y capacitación en entornos controlados.


Al tratarse de una solución personalizable, el
sistema puede adaptarse a requerimientos específicos del laboratorio, como el formato de
registro, la interfaz de usuario, los métodos de identificación y las condiciones de validación.

Por lo tanto, el presupuesto del proyecto respalda la viabilidad económica de la solución
propuesta y refuerza su principal aporte: desarrollar una estación ESD inteligente, de bajo
costo y adaptable, orientada a mejorar la trazabilidad y el cumplimiento de procedimientos
ESD en entornos de laboratorio.
